% Part: first-order-logic
% Chapter: natural-deduction
% Section: derivations

\documentclass[../../../include/open-logic-section]{subfiles}

\begin{document}

\iftag{FOL}
      {\olfileid{fol}{ntd}{der}}
      {\olfileid{pl}{ntd}{der}}

\olsection{\usetoken{P}{derivation}}

\begin{explain}
गृहीतक म्हणजे काय हे आपण सांगितले आहे आणि निगमन नियम दिले आहेत.
नैसर्गिक निगमनातील !!^{derivation}s यांपासून आगमनाने निर्माण होतात:
प्रत्येक !!{derivation} ही एकतर स्वतंत्र गृहीतक असते किंवा एक, दोन वा तीन
!!{derivation}s नंतर केलेले योग्य अनुमान असते.
\end{explain}

\begin{defn}[!!^{derivation}]
\Article{derivation} \emph{!!{derivation}} म्हणजे गृहीतके~$\Gamma$
यांपासून !!a{sentence}~$!A$ ची पुढील अटी पूर्ण करणारी !!{sentence}s ची
सांत वृक्षरचना होय:
\begin{enumerate}
\item वृक्षातील सर्वांत वरची !!{sentence}s ही एकतर $\Gamma$ मध्ये असतात
  किंवा वृक्षातील एखाद्या अनुमानाने !!{discharged} केलेली असतात.
\item वृक्षातील सर्वांत खालचे !!{sentence} हे~$!A$ असते.
\item वृक्षाच्या तळाशी असलेले वाक्य~$!A$ वगळता, वृक्षातील प्रत्येक
  !!{sentence} हे एखाद्या निगमन नियमाच्या योग्य उपयोजनाचे आधारविधान असते
  आणि त्या उपयोजनाचा निष्कर्ष वृक्षात त्या !!{sentence} च्या थेट खाली येतो.
\end{enumerate}
मग $!A$ हा त्या !!{derivation} चा \emph{निष्कर्ष} आणि $\Gamma$ ही तिची
!!{undischarged} गृहीतके आहेत, असे आपण म्हणतो.

जर~$\Gamma$ या गृहीतकांपासून $!A$ ची !!a{derivation} अस्तित्वात असेल, तर
$!A$ हे~$\Gamma$ पासून \emph{!!{derivable}} आहे असे म्हणतो; चिन्हांत:
$\Gamma \Proves !A$. प्रत्येक गृहीतक !!{discharged} केलेले असणारी $!A$
ची !!a{derivation} असेल, तर आपण~$\Proves !A$ असे लिहितो.
\end{defn}

\begin{ex}
प्रत्येक गृहीतक स्वतःच !!a{derivation} असते. म्हणून, उदा., एकटे $!A$ ही
!!a{derivation} आहे आणि एकटे $!B$ सुद्धा. यांवर, समजा,
$\Intro{\land}$ नियम लागू करून नवी !!{derivation} मिळवता येते:
\begin{prooftree}
\AxiomC{$!A$}
\AxiomC{$!B$}
\RightLabel{\Intro{\land}}
\BinaryInfC{$!A \land !B$}
\end{prooftree}
हे नियम सर्वसाधारण असतात: त्यांतील $!A$ आणि~$!B$ यांच्या जागी कोणतीही
!!{sentence}s, उदा., $!C$ आणि~$!D$, ठेवता येतात. मग निष्कर्ष
$!C \land !D$ असेल, आणि म्हणून
\begin{prooftree}
\AxiomC{$!C$}
\AxiomC{$!D$}
\RightLabel{\Intro{\land}}
\BinaryInfC{$!C \land !D$}
\end{prooftree}
ही योग्य !!{derivation} आहे. अर्थात, गृहीतकांची अदलाबदल करून $!D$ ला
$!A$ ची आणि $!C$ ला~$!B$ ची भूमिका देता येते. त्यामुळे,
\begin{prooftree}
\AxiomC{$!D$}
\AxiomC{$!C$}
\RightLabel{\Intro{\land}}
\BinaryInfC{$!D \land !C$}
\end{prooftree}
हीसुद्धा योग्य !!{derivation} आहे.

आता आपण दुसरा नियम, समजा, $\Intro{\lif}$ लागू करू शकतो. त्यामुळे सोपाधिक
विधानाचा निष्कर्ष काढता येतो आणि त्या विधानाच्या पूर्वांगाशी एकरूप असलेले
कोणतेही गृहीतक !!{discharge} करता येते. म्हणून पुढील दोन्ही
!!{derivation}s योग्य आहेत:
\begin{prooftree}
\AxiomC{$\Discharge{!C}{1}$}
\AxiomC{$!D$}
\RightLabel{\Intro{\land}}
\BinaryInfC{$!C \land !D$}
\DischargeRule{\Intro{\lif}}{1}
\UnaryInfC{$!C \lif (!C \land !D)$}
\DisplayProof\bottomAlignProof
\AxiomC{$!C$}
\AxiomC{$\Discharge{!D}{1}$}
\RightLabel{\Intro{\land}}
\BinaryInfC{$!C \land !D$}
\DischargeRule{\Intro{\lif}}{1}
\UnaryInfC{$!D \lif (!C \land !D)$}
\end{prooftree}
त्या अनुक्रमे $!D \Proves !C \lif (!C \land !D)$ आणि
$!C \Proves !D \lif (!C \land !D)$ हे दाखवतात.

गृहीतके मुक्त करणे ही परवानगी आहे, आवश्यकता नाही, हे आठवा: गृहीतके मुक्त
करणे बंधनकारक नसते. विशेषतः, !!{derivation} मध्ये ती गृहीतके नसली तरीही
नियम लागू करता येतो. उदाहरणार्थ, !!{discharged} करण्यासाठी~$!A$ हे गृहीतक
नसले तरी पुढील उपयोजन अनुज्ञात आहे:
\begin{prooftree}
  \AxiomC{$!B$}
  \DischargeRule{\Intro{\lif}}{1}
  \UnaryInfC{$!A \lif !B$}
\end{prooftree}
\end{ex}

\end{document}
